\subsection{主覆叠}

定义 2. —— 设 B 为拓扑空间，G 为群。给 G 赋予离散拓扑。以 B 为基、以 G 为群的主纤维丛称为 B 的以 G 为群的主覆叠。

这一术语是合理的，因为这样的 B-空间是 B 的一个覆叠。

非空拓扑空间上的主 G-丛具有次数，等于 Card G。

命题 4. —— 设 B 为拓扑空间，(E, p) 为 B-空间，G 为从右边作用于 E 的离散拓扑群。下列断言等价：

\begin{enumerate}
\def\labelenumi{(\roman{enumi})}
\item
  B-空间 E 是以 G 为群的主覆叠；
\item
  对每个 \(g \in G\) 和每个 \(x \in E\) ，有 \(p(x \cdot g) = p(x)\) ，映射 p 诱导从 E/G 到 B 上的同胚，且映射 \(\theta: (x, g) \mapsto (x, x \cdot g)\) 是从 E × G 到 \(E \times_{B} E\) 上的同胚；
\item
  群 G 连续且自由地作用于 E，映射 p 是平展的，且其纤维是 G 的轨道。
\end{enumerate}

(i)⇒(ii)：这由 I, 第 91 页以及 I, 第 94 页命题 3 得出。(ii)⇒(iii)：在假设 (ii) 之下，G 的作用法则是连续的，且映射 p 是满射且开映射（TG, III, 第 10 页, 引理 2）。设 \(x \in E\) ；映射 \(\theta\) 诱导从 \(\{x\} \times G\) 到 \(\{x\} \times p^{-1}(p(x))\) 上的双射，故群 G 自由作用，且其轨道是 p 的纤维。若以 e 表示 G 的单位元，则 \(E \times_{B} E\) 的对角线是 \(E \times \{e\}\) 在同胚 \(\theta\) 下的像。由于群 G 是离散的，集合 \(\{e\}\) 在 G 中是开集，从而对角线 \(\Delta_{E}\) 在 \(E \times_{B} E\) 中是开集。由 I, 第 31 页命题 7，映射 p 是平展的。

(iii)⇒(i)：由于 p 的每个纤维都是 G 在 E 上作用的一条轨道，映射 p 是满射。因为 p 是平展的，对 B 的每个点 b，存在 b 的开邻域 U 以及 p 在 U 上的连续截面 s。集合 \(s(\mathrm{U})\) 在 E 中是开集，且 s 诱导从 U 到 \(s(\mathrm{U})\) 上的同胚（I, 第 30 页, 推论 3）。对每个 \(g \in G\)，集合 \(s(\mathrm{U}) \cdot g\) 在 E 中是开集，且这些集合对一切 \(g \in G\) 的并等于 \(p^{-1}(\mathrm{U})\)。若 g 与 \(g'\) 是 G 的两个不同元素，则集合 \(s(\mathrm{U}) \cdot g\) 与 \(s(\mathrm{U}) \cdot g'\) 不相交，因为 G 自由地作用于 E。映射 \(f: U \times G \to p^{-1}(\mathrm{U})\) 由 \(f(u, g) = s(u) \cdot g\) 所定义，对 \((u, g) \in U \times G\) 是连续的，因而是从 \(U \times G\) 到 \(p^{-1}(\mathrm{U})\) 上的与 G 的右作用相容的同胚。按定义，E 因而是 B 的以 G 为群的主覆叠。

例. —— 1) 设 E 为拓扑空间，G 为从右边连续且自由地作用于 E 的离散拓扑群。为使典范映射 \(p \colon \mathrm{E} \to \mathrm{E} / \mathrm{G}\) 使 E 成为 E/G 的主覆叠，充要条件是该映射是平展的（命题 4 的条件 (iii)）。例如，假设存在拓扑空间 X 以及与 G 在 E 上的作用相容的平展映射 \(q \colon \mathrm{E} \to \mathrm{X}\) ；则 E 是 E/G 的主覆叠。事实上，以 \(q' \colon \mathrm{E} / \mathrm{G} \to \mathrm{X}\) 表示由 \(q\) 诱导的映射；我们有 \(q = q' \circ p\)，因此由 I, 第 29 页命题 6 d)，\(p\) 是平展的。

\begin{enumerate}
\def\labelenumi{\arabic{enumi})}
\setcounter{enumi}{1}
\tightlist
\item
  设 \(q \colon \mathrm{E} \to \mathrm{X}\) 为分离平展态射。群 \(\operatorname{Aut}_{\mathrm{X}}(\mathrm{E})\) 赋予离散拓扑后，从左边连续地作用于 E。若空间 E 连通，则此作用是自由的（I, 第 34 页, 推论 2）。以 G 表示群 \(\operatorname{Aut}_{\mathrm{X}}(\mathrm{E})^{\circ}\) ，并给 E 赋予与 \(\operatorname{Aut}_{\mathrm{X}}(\mathrm{E})\) 的作用相反的右作用。由前例，赋予典范映射 \(p \colon \mathrm{E} \to \mathrm{E} / \mathrm{G}\) 的空间 E 是以 G 为群的主覆叠。
\end{enumerate}

